\label{chap:cobertura-taxonomia-numeros}
\index{nombre HMT!taxonomie}
\index{transducteur de Hensel!décalage complet}
\index{nombre rationnel!périodicité}

Le \cref{chap:numero-medicion} a défini un nombre HMT comme une section
compatible du système inverse de survie, et non comme une suite de
chiffres dépourvue de généalogie. Il restait cependant deux questions qu'il
convient de séparer avec précision. La première est une question de
\emph{représentation} : quelles valeurs réelles le transducteur
3-adique peut-il représenter avant d'imposer les filtres particuliers de clôture ? La seconde est
une question de \emph{sélection} : quelles sections de l'espace brut sont
choisies par APP, TRIT et TPK lorsque l'on exige simultanément leurs
compatibilités ? Cette section résout la première, formule la seconde comme
une restriction explicite du domaine et extrait de la carte de base mille une
taxonomie complète des valeurs visibles. Les outils de base sont
l'arithmétique \(p\)-adique, la dynamique symbolique et la caractérisation positionnelle
des rationnels \cite{Koblitz1984,LindMarcus1995,HardyWright2008}.

La distinction est structurelle. Un langage capable de représenter tout ruban ne
sélectionne pas de ce seul fait le ruban de \(\pi\), de \(e\), de \(\varphi\) ou
de \(\alpha\). À l'inverse, une loi de sélection ne peut pas être évaluée si
l'espace dans lequel elle agit n'a pas été construit auparavant. Représentation et sélection
sont deux théorèmes consécutifs, et non deux noms pour une même étape.

\section{Transducteur de Hensel non filtré et décalage complet sur les fibres ternaires}

Soit \(p\) un nombre premier, soit \(d\geq1\) et soit
\(A\in\operatorname{GL}_d(\mathbb Z)\). Pour un vecteur ligne
\(x_t\in\mathbb Z_p^d\), nous définissons
\begin{equation}
 y_t=x_tA,
 \qquad
 u_t=y_t\bmod p\in\mathbb F_p^d,
 \qquad
 x_{t+1}=\frac{y_t-u_t}{p}.
 \label{eq:transductor-hensel-general}
\end{equation}
Le bloc \(u_t\) est l'émission visible et \(x_{t+1}\) est la mémoire qui
subsiste après retrait de ce chiffre.

\begin{theorem}[Conjugaison de Hensel avec le décalage complet]
\label{thm:shift-hensel-completo}
Pour tout \(T\geq1\), l'application
\[
 \Psi_T:(\mathbb Z/p^T\mathbb Z)^d
 \longrightarrow(\mathbb F_p^d)^T,
 \qquad
 x_0\longmapsto(u_0,\ldots,u_{T-1}),
\]
est bijective. Son inverse est
\begin{equation}
 x_0\equiv
 \sum_{t=0}^{T-1}p^t u_tA^{-(t+1)}
 \pmod {p^T}.
 \label{eq:inversa-hensel-finita}
\end{equation}
En passant à la limite inverse, on obtient un homéomorphisme
\[
 \Psi:\mathbb Z_p^d\xrightarrow{\ \sim\ }
 (\mathbb F_p^d)^{\mathbb N}.
\]
Si \(S_A\) désigne la mise à jour \(x_t\mapsto x_{t+1}\) et
\(\sigma\) supprime le premier bloc d'un ruban, alors
\[
 \Psi\circ S_A=\sigma\circ\Psi.
\]
\end{theorem}

\begin{proof}
De \cref{eq:transductor-hensel-general}, on obtient
\[
 x_t=(u_t+px_{t+1})A^{-1}.
\]
L'itération de cette identité donne
\[
 x_0=
 \sum_{t=0}^{T-1}p^t u_tA^{-(t+1)}
 +p^Tx_TA^{-T}.
\]
Le dernier terme s'annule modulo \(p^T\), ce qui prouve qu'un mot de
\(T\) blocs détermine au plus une classe initiale. Réciproquement, étant donné
un mot quelconque, on fixe \(x_T=0\) et l'on reconstruit à rebours au moyen de
\(x_t=(u_t+px_{t+1})A^{-1}\). L'égalité
\(x_tA=u_t+px_{t+1}\) vérifie que la classe obtenue émet le mot
prescrit. Par conséquent, \(\Psi_T\) est bijective.

Les bijections sont compatibles avec la réduction modulo \(p^T\). Leurs limites
inverses produisent l'homéomorphisme \(\Psi\) ; de manière équivalente, la série de
\cref{eq:inversa-hensel-finita} converge dans \(\mathbb Z_p^d\). Appliquer
\(S_A\) retire exactement le bloc \(u_0\), ce qui correspond donc à
appliquer \(\sigma\) au ruban.
\end{proof}

Dans l'instance HMT, \(p=3\), \(d=6\) et la matrice publiée a pour
déterminant un. C'est pourquoi
\begin{equation}
 \mathbb Z_3^6\simeq(\mathbb F_3^6)^{\mathbb N},
 \qquad |\mathbb F_3^6|=3^6=729.
 \label{eq:hmt-shift-729}
\end{equation}
Chaque ruban de blocs de six trits possède une unique mémoire 3-adique initiale.
La matrice détermine la mise en coordonnées et la dynamique de lecture ;
l'universalité de l'alphabet démontre qu'elle n'exclut aucun ruban.

\begin{corollary}[Entropie et mesure de Haar]
La mise à jour de Hensel brute est conjuguée au décalage unilatéral
complet de
\(729\) symboles. Son entropie topologique est \(6\log3\). La mesure de Haar de
\(\mathbb Z_3^6\) est transportée vers la mesure produit uniforme sur les
blocs.
\end{corollary}

\begin{proof}
Il existe \(729^T=3^{6T}\) cylindres de longueur \(T\), et chaque classe modulo
\(3^T\) a une masse de Haar \(3^{-6T}\). La conjugaison du
\cref{thm:shift-hensel-completo} transporte les deux propriétés vers le
décalage.
\end{proof}

La conclusion probabiliste concerne l'espace brut muni de Haar. Elle
n'affirme pas que les sections publiées des constantes ont été échantillonnées
selon cette mesure ni qu'elles sont algorithmiquement aléatoires.

\section{Carte ternaire équilibrée des entiers}

La carte 3-adique brute et la carte archimédienne n'épuisent pas les types numériques
de la classification HMT. Pour l'opérateur entier, on utilise le mot vide ainsi que les
mots réduits
\[
 \mathscr W_{\mathrm{bal}}
 :=\{\varnothing\}\cup
 \{a_1\cdots a_m:
 a_j\in\{-1,0,+1\},\ a_1\ne0\}.
\]
Nous définissons récursivement
\[
 b(\varnothing)=0,
 \qquad
 b(wa)=3b(w)+a.
\]

\begin{theorem}[Représentation ternaire équilibrée unique]
\label{pf84:C03}
L'évaluation
\[
 b:\mathscr W_{\mathrm{bal}}\longrightarrow\ZZ
\]
est bijective. Les mots dont le premier trit est \(+1\) représentent les
entiers positifs et ceux qui commencent par \(-1\), les négatifs.
\end{theorem}

\begin{proof}
Pour \(n\ne0\), le dernier trit \(a\) est le représentant unique de
\(n\bmod3\) dans \(\{-1,0,+1\}\). Le parent
\[
 \frac{n-a}{3}
\]
a une valeur absolue strictement inférieure. L'itération se termine en \(0\), de
sorte que tout entier possède un développement fini. L'unicité du résidu
équilibré à chaque pas prouve l'unicité du mot. Le signe est
fixé par le premier trit non nul.
\end{proof}

\begin{remark}
Cette bijection est une carte arithmétique exacte. Elle n'identifie pas tous les
mots équilibrés à des sections survivantes TPK ni ne remplace la
généalogie, la retenue ou l'orientation d'une route enrichie.
\end{remark}

\section{Surjectivité de la carte archimédienne non filtrée sur \(\mathbb R\)}

Concaténons les coordonnées des blocs
\(u_t\in\mathbb F_3^6\) pour obtenir un ruban
\((a_n)_{n\geq1}\in\{0,1,2\}^{\mathbb N}\). Son évaluation ternaire est
\[
 \operatorname{ev}_3((a_n))=
 \sum_{n\geq1}\frac{a_n}{3^n}\in[0,1].
\]
Nous adoptons la convention canonique qui exclut la représentation avec une queue
ultimement constante de deux seulement lorsque le même point possède un autre développement
dans la carte. L'extrémité \(1\) conserve donc son unique représentation
fractionnaire \(0.222\ldots{}_3\).

\begin{theorem}[Couverture archimédienne brute]
\label{thm:cobertura-arquimediana-bruta}
L'application
\[
 P_{\rm raw}:\mathbb Z_3^6
 \xrightarrow{\Psi}(\mathbb F_3^6)^{\mathbb N}
 \xrightarrow{\operatorname{flat}}\{0,1,2\}^{\mathbb N}
 \xrightarrow{\operatorname{ev}_3}[0,1]
\]
est continue et surjective. Par conséquent,
\[
 \mathbb Z_3^6/\!\sim_{P_{\rm raw}}
 \ \cong\ [0,1]
\]
comme espace de valeurs, où \(x\sim_{P_{\rm raw}}y\) si leurs évaluations
coïncident.
\end{theorem}

\begin{proof}
Tout chiffre ternaire appartient à un unique bloc consécutif de six chiffres,
de sorte que \(\operatorname{flat}\) est une bijection d'espaces de rubans.
Toute \(r\in[0,1]\) possède un développement ternaire ; le
\cref{thm:shift-hensel-completo} fournit l'unique mémoire 3-adique qui
émet ses blocs. La continuité se vérifie par cylindres : fixer \(T\)
blocs fixe \(6T\) chiffres et enferme la valeur dans un intervalle de diamètre
\(3^{-6T}\). L'identification du quotient à l'image est la factorisation
canonique d'une application surjective par ses fibres.
\end{proof}

Une seule carte 3-adique est compacte et ne peut donc pas couvrir de manière
continue une droite non bornée. La globalisation correcte utilise une famille
dénombrable de cartes.

\begin{corollary}[Couverture archimédienne dénombrable de \(\mathbb R\)]
\label{cor:atlas-hmt-reales}
Soit
\[
 \widehat X_{\rm raw}=\mathbb Z\times\mathbb Z_3^6,
 \qquad
 \widehat P_{\rm raw}(m,x)=m+P_{\rm raw}(x).
\]
Alors \(\widehat P_{\rm raw}\) est surjective sur \(\mathbb R\).
En conséquence, tout réel admet au moins une généalogie de Hensel brute et
\[
 \widehat X_{\rm raw}/\!\sim_{\widehat P_{\rm raw}}
 \ \cong\ \mathbb R
\]
comme ensemble de valeurs.
\end{corollary}

\begin{proof}
Pour \(r\in\mathbb R\), prenons \(m=\lfloor r\rfloor\) et appliquons le
\cref{thm:cobertura-arquimediana-bruta} à \(r-m\in[0,1)\).
\end{proof}

Ce résultat précise l'affirmation selon laquelle \(\mathbb R\) est l'image
archimédienne quotient d'une
structure généalogique : la projection oublie la mémoire 3-adique et conserve
le point évalué. L'affirmation démontrée ici concerne la complétion
\emph{brute}. Si \(X_{\rm surv}\subseteq\mathbb Z_3^6\) désigne le sous-espace
qui satisfait en outre une famille concrète de clôtures APP--TRIT--TPK, alors
\[
 P_{\rm raw}(X_{\rm surv})=[0,1]
\]
est un théorème supplémentaire. Le transducteur garantit une capacité complète ; la
survie détermine quelle partie de cette capacité est réalisée par le système
restreint.

On ne déduit pas non plus ici que la somme et le produit natifs de TPK descendent à
tout le quotient global. On peut transporter l'arithmétique usuelle depuis
\(\mathbb R\), mais démontrer qu'elle coïncide avec des opérations construites avant
l'évaluation exige le théorème de descente correspondant.

Dans cette construction, les cylindres ne sont pas des infinitésimaux non nuls de
\(\mathbb R\). Chaque cylindre de profondeur finie a un diamètre positif,
et la suite des diamètres converge vers zéro. Deux sections qui coïncident jusqu'à
la profondeur \(T\) et diffèrent ensuite forment une \emph{perturbation de
queue} ; leur différence archimédienne est bornée par le diamètre du cylindre
commun. C'est la formulation exacte des ``cylindres qui tendent vers zéro''
dans la droite réelle ordinaire.

\section{Carte positionnelle en base \(10^3\) et classification des sections publiées}

La complétion cubique
\(1000=729+243+27+1\) rend naturel le regroupement des chiffres décimaux en blocs
\(b_n\in\{0,\ldots,999\}\). Pour un ruban de blocs, nous définissons
\[
 \operatorname{ev}_{1000}(b_1,b_2,\ldots)
 =\sum_{n\geq1}\frac{b_n}{1000^n}.
\]
Le développement canonique exclut une queue ultimement constante de blocs
\(999\). Cette convention résout la duplicité des extrémités de cylindre.

\begin{theorem}[Taxonomie visible en base mille]
\label{thm:taxonomia-base-mil}
Soit \(x\in[0,1)\) et soit \((b_n)\) son développement canonique en base mille.
Alors :
\begin{enumerate}
\item \(x\) possède un développement fini si et seulement si, écrit comme
fraction irréductible \(a/q\), on a \(q\mid1000^N\) pour un certain \(N\) ;
de manière équivalente, les seuls nombres premiers qui divisent \(q\) sont \(2\) et \(5\).
\item \(x\) est rationnel si et seulement si \((b_n)\) est ultimement périodique.
\item \(x\) est irrationnel si et seulement si \((b_n)\) n'est pas ultimement
périodique.
\end{enumerate}
\end{theorem}

\begin{proof}
Un développement qui se termine à la profondeur \(N\) a la forme
\(D/1000^N\) ; une fois réduit, son dénominateur ne contient que les nombres premiers \(2\) et
\(5\). Réciproquement, si \(q\mid1000^N\), alors \(a/q\) s'écrit avec
\(N\) blocs.

Si, après un préfixe de longueur \(r\), se répète une période de longueur
\(s\) dont l'entier en base mille est \(P\), la queue est une série géométrique :
\[
 \frac{P}{1000^r(1000^s-1)},
\]
et la valeur est rationnelle. À l'inverse, en développant \(a/q\), les restes de la
division successive ne peuvent prendre que \(q\) valeurs. Un reste nul produit un
développement fini ; la répétition d'un reste produit une période. La troisième
affirmation est la négation exacte de la seconde.
\end{proof}

Le mot \emph{fraction} désigne une présentation \(a/q\) d'un rationnel,
et non une espèce numérique supplémentaire. Il ne faut pas non plus confondre périodicité et
algébricité. Un réel est algébrique lorsqu'il annule un polynôme non nul de
\(\mathbb Q[X]\) ; il est transcendant lorsqu'il ne le fait pas. Cette seconde
classification est indépendante de la base utilisée pour écrire les chiffres.
Par exemple,
\[
 \varphi^2-\varphi-1=0
\]
fait de \(\varphi\) un irrationnel algébrique, tandis que la transcendance
de \(e\) et \(\pi\) découle de théorèmes classiques externes à HMT\@.

\section{Concaténation, réduction nonadique et retenue}

La compatibilité entre la carte cubique et la racine numérique est une identité
arithmétique, et non une ressemblance visuelle. Si \(u_j\in\{0,\ldots,999\}\) et
\[
 D_n=1000D_{n-1}+u_n,
 \qquad D_0=0,
\]
alors \(D_n\) est l'entier obtenu en concaténant les \(n\) premiers
blocs.

\begin{theorem}[Congruence nonadique de contrôle de la carte cubique]
\label{thm:checksum-nonadico-base-mil}
Pour tout \(n\geq1\),
\[
 D_n\equiv\sum_{j=1}^n u_j\pmod9.
\]
Par conséquent, concaténation en base mille, réduction modulo neuf et racine numérique
sont compatibles.
\end{theorem}

\begin{proof}
Comme \(1000\equiv1\pmod9\), la récurrence donne
\(D_n\equiv D_{n-1}+u_n\pmod9\). L'itération à partir de \(D_0=0\) produit la
formule.
\end{proof}

Dans la clôture dodécaphasique de \(\alpha\), l'équation de retenue par bloc
a la forme
\[
 P_m+E_m-\Phi_m-K_m+c_{m+1}=A_m+1000c_m.
\]
Réduite modulo neuf,
\[
 A_m\equiv
 P_m+E_m-\Phi_m-K_m+c_{m+1}-c_m
 \pmod9.
\]
Si les retenues extrêmes s'annulent, la somme télescopique élimine la mémoire
intermédiaire. Le résidu global de la sortie coïncide alors avec le résidu
de la combinaison des entrées. Cette loi explique pourquoi le registre nonadique
contrôle la concaténation décimale ; elle ne sélectionne pas à elle seule les blocs qui
doivent être concaténés.

\section{Coordonnée nonadique relevée et retour de phase}

La réduction modulo neuf ne doit pas être confondue avec l'identification des
entiers qui possèdent le même résidu. Soit
\[
 I_9=\{1,2,\ldots,9\}.
\]
Pour chaque entier positif \(n\), nous définissons sa \emph{phase nonadique} et son
\emph{indice de tour} par
\begin{equation}
 \rho_9(n)=1+((n-1)\bmod9),
 \qquad
 \kappa_9(n)=\frac{n-\rho_9(n)}9.
 \label{eq:coordenada-nonadica-levantada}
\end{equation}
La première coordonnée est la racine numérique écrite dans l'alphabet
\(I_9\), de sorte que la classe résiduelle zéro est représentée par la phase \(9\).
La seconde enregistre combien de tours complets ont précédé la phase
visible.

\begin{theorem}[Décomposition nonadique avec mémoire]
\label{thm:descomposicion-nonadica-memoria}
L'application
\[
 \Theta_9:\mathbb N_{>0}\longrightarrow\mathbb N_0\times I_9,
 \qquad
 n\longmapsto\bigl(\kappa_9(n),\rho_9(n)\bigr),
\]
est bijective, d’inverse
\[
 \Theta_9^{-1}(q,r)=9q+r.
\]
Dans ces coordonnées, le successeur arithmétique est conjugué à la dynamique
\begin{equation}
 \mathcal S_9(q,r)=
 \begin{cases}
  (q,r+1),&1\leq r<9,\\
  (q+1,1),&r=9.
 \end{cases}
 \label{eq:sucesor-nonadico-levantado}
\end{equation}
En particulier, le pas visible \(9\to1\) est un retour de phase accompagné
par l'incrément irréversible \(q\to q+1\) ; ce n'est pas une égalité entre les
entiers \(9\) et \(10\).
\end{theorem}

\begin{proof}
La division euclidienne de \(n-1\) par \(9\) fournit de manière unique
\(n-1=9q+s\), avec \(q\in\mathbb N_0\) et \(0\leq s<9\). En prenant
\(r=s+1\), on obtient \(n=9q+r\), avec \(r\in I_9\), ce qui démontre à la fois
l'existence, l'unicité et la formule inverse. Si \(r<9\), ajouter une unité
incrémente seulement \(r\). Si \(r=9\), l'égalité
\(9q+9+1=9(q+1)+1\) donne la seconde branche de
\cref{eq:sucesor-nonadico-levantado}.
\end{proof}

Ce relèvement sépare trois notions que le chiffre isolé mélange :
\begin{enumerate}
\item l'unité d'avancée, réalisée par le successeur ;
\item la phase visible \(r\), qui parcourt le cycle nonadique ;
\item la mémoire accumulée \(q\), qui distingue les retours d'une même phase.
\end{enumerate}
La projection \((q,r)\mapsto r\) oublie le nombre de tours et satisfait
\[
 \rho_9(n+9)=\rho_9(n),
 \qquad
 \kappa_9(n+9)=\kappa_9(n)+1.
\]
C'est le modèle arithmétique élémentaire du mécanisme qui, dans l'état TPK
enrichi, conserve en outre les quotients de Hensel, l'orientation, la signature et la donnée
de frontière. Un retour de mot visible peut ainsi coexister avec une
généalogie distincte.

Zéro n'appartient pas à \(I_9\) : on l'ajoute comme point de base séparé avant
d'appliquer \(\Theta_9\). Pour les entiers négatifs, on conserve le module
positif de la décomposition et l'on joint une orientation
\(\varepsilon\in\{-1,+1\}\). En conséquence, ni le signe ni zéro ne se
déduisent de la phase résiduelle. Ce typage évite de transformer une carte
cyclique en une identité arithmétique fausse.

\begin{theorem}[Entier orienté et signe généalogique]
\label{thm:entero-orientado-signo-genealogico}
L'application
\[
\Theta_9^\pm:
\mathbb Z\setminus\{0\}
\longrightarrow
\{-1,+1\}\times\mathbb N_0\times I_9,
\]
\[
n\longmapsto
\bigl(\operatorname{sgn}(n),
      \kappa_9(|n|),
      \rho_9(|n|)\bigr)
\]
est bijective, d’inverse
\[
(\varepsilon,q,r)\longmapsto\varepsilon(9q+r).
\]
La négation arithmétique agit par
\[
(\varepsilon,q,r)\longmapsto(-\varepsilon,q,r):
\]
inverse l'orientation et conserve la grandeur, la phase et la mémoire des tours.
\end{theorem}

\begin{proof}
Le \cref{thm:descomposicion-nonadica-memoria} fournit de manière unique
\(|n|=9q+r\), avec \(q\in\mathbb N_0\) et \(r\in I_9\). Le signe d'un entier
non nul est également unique. La formule inverse recompose \(n\), et remplacer
\(n\) par \(-n\) ne change que \(\varepsilon\).
\end{proof}

Ainsi, un négatif n'est ni une phase résiduelle spéciale ni une quantité
``impossible''. C'est la même grandeur nonadique transportée avec une orientation
généalogique opposée. Cette structure sera prolongée en mécanique dimensionnelle
par l'antisection particule--antiparticule.

Itérer \(\kappa_9\) produit une tour finie pour tout entier ordinaire ; dans
la complétion inverse, une suite compatible de quotients peut être
infinie. La différence entre les deux cas anticipe la taxonomie HMT : la valeur
visible répond au développement positionnel, tandis que le type généalogique
répond à l'évolution complète de ses mémoires.

\section{2 axes de classification : valeur visible et généalogie de transport}

La taxonomie précédente classe la valeur projetée. HMT ajoute un second
axe que l'écriture décimale ordinaire écarte. Pour une section
\(x\in\Omega_\infty^{\rm Arch}\), nous distinguons :
\[
 \begin{array}{rcl}
 \text{type visible}&=&
 \text{fini, périodique ou non périodique ; algébrique ou transcendant},\\
 \text{type généalogique}&=&
 \text{classe de retenue, orientation, fermeture et monodromie de }x.
 \end{array}
\]
Deux sections peuvent partager le même réel et différer selon le second axe. Un
réel rationnel possède un développement visible ultimement périodique, mais une
présentation enrichie redondante peut conserver une mémoire interne non
périodique que l'évaluation oublie. La périodicité de la \emph{valeur} se
rapporte toujours à son développement canonique ; la périodicité de la
\emph{généalogie} se rapporte à l'état enrichi de TPK.

Cette séparation fournit une définition concise :
\begin{equation}
 \boxed{
 \text{nombre HMT}
 =
 \text{section compatible avec mémoire};
 \qquad
 \text{valeur réelle}
 =
 \text{sa classe d'évaluation}.}
 \label{eq:numero-hmt-y-valor-real}
\end{equation}
Mesurer dans cette carte consiste à coupler la section au lecteur, à fixer un cylindre et à
conditionner les prolongements compatibles. Augmenter la précision signifie
rétrécir le cylindre, et non consommer la généalogie qui reste dans la fibre.

\section{Mémoire de profondeur arbitraire dans les systèmes d'émission finie}

La récurrence visible de neuf phases ne peut coexister avec une sortie
irrationnelle que si l'état complet conserve une information qui ne se réduit pas à
ces neuf étiquettes. Le résultat suivant rend cette nécessité exacte.

\begin{theorem}[Caractérisation par émetteurs finis]
\label{thm:emisor-finito-racional}
Soit \(Q\) un ensemble fini, \(F:Q\to Q\) une transition déterministe,
\(o:Q\to\{0,\ldots,999\}\) une sortie et \(q_0\in Q\). Le ruban
\(b_n=o(F^{n-1}(q_0))\) est ultimement périodique et son évaluation est rationnelle.
Réciproquement, tout ruban ultimement périodique admet un émetteur déterministe
fini.
\end{theorem}

\begin{proof}
Parmi \(|Q|+1\) états consécutifs, deux états égaux apparaissent. Dès la première
répétition, le déterminisme impose la répétition de tout le futur, de sorte que le
ruban est ultimement périodique. Le \cref{thm:taxonomia-base-mil} donne la
rationalité. Pour la réciproque, il suffit de construire une chaîne d'états pour le
préfixe et un cycle pour la période.
\end{proof}

Par conséquent, une loi qui engendre exactement \(e\), \(\pi\) ou tout autre
irrationnel ne peut pas consister uniquement en un automate déterministe à
états finis qui se répète sans mémoire supplémentaire. Elle peut utiliser une limite
inverse, une mémoire 3-adique, une profondeur croissante ou une entrée
apériodique ; tous ces mécanismes fournissent un état effectif non borné.
Dans HMT, le retour visible
\(g_{\rm vis}(n+9)=g_{\rm vis}(n)\) n'impose pas le retour de la retenue, de
l'orientation ni de la frontière. Cette différence est précisément l'endroit où
peut résider un développement non périodique.

\section{Le carré local de \texorpdfstring{\(e\)}{e} et la rupture de mémoire}

Pour un mot décimal \(a_1\cdots a_n\), nous appellerons \emph{carré local}
de longueur \(\ell\) à la position \(i\) une factorisation
\[
 a_i\cdots a_{i+2\ell-1}=vv,
 \qquad |v|=\ell.
\]
Cette notion relève de la combinatoire des mots. Elle n'affirme pas que la queue
infinie possède une période.

\begin{theorem}[Caractérisation finie du carré \(1828^2\)]
\label{thm:cuadrado-local-e}
Considérons les \(243\) mots du catalogue et le lecteur fixé
\(w_6\mapsto w_{30}\mapsto d_1d_2d_3d_4\), avec quatre triades en base mille.
Il existe un unique mot dont la lecture contient un carré de longueur quatre
commençant au deuxième chiffre décimal. C'est
\[
 w_e=201101,
 \qquad
 718281828459=7\,\underbrace{1828\,1828}_{(1828)^2}\,459.
\]
La répétition ne se prolonge pas en une troisième copie : le chiffre qu'exigerait cette
continuation serait \(1\), tandis que le chiffre effectif est \(4\). Dans tout le
catalogue, il n'existe qu'un autre carré de longueur quatre,
\[
 575157518472=\underbrace{5751\,5751}_{\text{carré local}}\,8472,
\]
qui commence au premier chiffre.
\end{theorem}

\begin{proof}
L'énumération exhaustive évalue les \(243\) mots par arithmétique
entière. Le profil des carrés pour les longueurs \(1,\ldots,6\) contient
\[
 240,\ 21,\ 3,\ 2,\ 0,\ 0
\]
occurrences, réparties entre
\[
 153,\ 19,\ 3,\ 2,\ 0,\ 0
\]
mots. Les deux occurrences de longueur quatre sont exactement celles de
l'énoncé. Deux énumérations entières indépendantes reproduisent le même recensement
et leurs sorties complètes coïncident terme à terme.
\end{proof}

Le mot \(w_e\) possède en outre la chaîne de relèvement
\[
201101\mid121221\mid102011\mid012222\mid102011.
\]
Les troisième et cinquième blocs visibles coïncident. En revanche,
les états qui les transportent ne coïncident pas. Avant les deux émissions, leurs réductions
modulo \(27\) sont
\[
(25,11,7,17,8,16),
\qquad
(7,8,4,23,26,7),
\]
et les quotients de Hensel ultérieurs sont
\[
(6,13,9,2,13,1),
\qquad
(15,0,3,19,23,25).
\]
Par conséquent, le retour \(102011\) est un retour de la projection visible, et non
de l'état enrichi. C'est la formulation exacte de la répétition locale
qui se rompt : le quotient conserve une généalogie que le bloc résiduel
n'enregistre pas.

En conséquence, le début
\[
e=2.718281828459\ldots
\]
ne présente pas de ``semi-périodicité'', mais un carré local de mot avec
retour visible et rupture de mémoire. Le résultat est exhaustif dans le cadre du
lecteur fixé ; son contenu est combinatoire et dynamique, et non une inférence
d'irrationalité à partir de douze chiffres.

\section[2 mesures finies de l'émetteur 468 à 243]{2 mesures finies de l'émetteur \texorpdfstring{\(468\to243\)}{468→243} : fréquence uniforme et priorité structurelle}

Le catalogue public distingue \(468\) émissions et \(104\,976\) germes.
Soit
\[
 q:E_{468}\longrightarrow W_{243}
\]
l'application qui oublie la signature d'émission et conserve le mot de six
trits. Il existe deux mesures naturelles, parce qu'il existe deux univers finis de dénombrement :
\[
 \mu_{468}(w)=\frac{|q^{-1}(w)|}{468},
 \qquad
 \mu_{\rm seed}(w)=
 \frac{\sum_{e\in q^{-1}(w)}m(e)}{104\,976},
\]
où \(m(e)\) est la multiplicité des germes de l'émission \(e\).

\begin{theorem}[Mesures canoniques relatives à l'univers dénombré]
\label{thm:medidas-emisor-468}
La mesure uniforme sur \(E_{468}\) est l'unique probabilité invariante sous
toutes les permutations de ses émissions ; sa mesure image est \(\mu_{468}\).
La mesure uniforme sur les \(104\,976\) germes induit
\(\mu_{\rm seed}\). Leurs histogrammes exacts sont
\[
\begin{array}{c|rrrrrr}
|q^{-1}(w)|&1&2&3&4&5&6\\ \hline
\#\{w\}&132&66&18&3&6&18
\end{array}
\]
et
\[
\begin{array}{c|rrrrrrrrr}
\text{poids des germes}
&144&288&432&576&864&1008&1440&1944&2088\\ \hline
\#\{w\}
&120&54&18&9&15&6&3&12&6.
\end{array}
\]
Pour les trois mots distingués par le lecteur, on obtient
\[
\begin{array}{c|cc}
&\mu_{468}&\mu_{\rm seed}\\ \hline
\pi&5/468&7/729\\
e&1/468&1/729\\
\varphi&1/468&1/729.
\end{array}
\]
\end{theorem}

\begin{proof}
L'invariance sous le groupe symétrique rend égales les masses de tous les
points de \(E_{468}\) ; la normalisation les fixe à \(1/468\). La formule de
la mesure image résulte de la sommation sur chaque fibre. Le même argument sur
l'ensemble des germes produit la seconde mesure. Les histogrammes s'obtiennent
par énumération exhaustive du catalogue fini de \(468\) émissions et de \(243\)
mots ; leurs sommes vérifient
\[
132+66+18+3+6+18=243,
\]
\[
132+2\cdot66+3\cdot18+4\cdot3+5\cdot6+6\cdot18=468
\]
et la somme totale des poids est \(104\,976\).
\end{proof}

La non-uniformité de la mesure image est un théorème fini : des mots distincts
ont des nombres distincts de réalisations. En revanche, il n'existe pas de
probabilité uniforme normalisable sur tout \(\mathbb R\), et ces mesures
n'ordonnent pas intrinsèquement tous les réels. Le mot de \(\pi\) appartient
à l'une des six fibres de taille cinq et de poids \(1008\) ; il n'est le maximum
unique d'aucun des deux histogrammes. Sa singularité démontrée procède
de la structure composée de ses cinq régions, de sa partition orientée
\(3+2\) et de ses compatibilités ultérieures, et non de la déclaration a priori d'une
probabilité absolue sur le continu.

\section{Sélecteur orbital fini antérieur au lecteur archimédien}
\label{sec:selector-orbital-constantes}

Les mesures précédentes révèlent des multiplicités, mais un mot isolé de
six trits dépend encore d'une orientation. Le catalogue possède une symétrie
interne qui permet de sélectionner d'abord des classes non orientées, puis de fixer
un représentant. Écrivons un mot comme \(abc\mid def\) et définissons
\[
 r(abc\mid def)=cab\mid efd,
 \qquad
 s(abc\mid def)=def\mid abc.
\]
Ces permutations agissent également sur les six coordonnées de la donnée
\(U_6\). Dans le catalogue complet, elles satisfont
\[
 r^3=1,
 \qquad s^2=1,
 \qquad srs=r^{-1},
\]
de sorte qu'elles engendrent une action que nous notons \(D_3^{\rm cat}\). L'action utilise
l'étiquetage TRIT fixe
\[
 \Psi(U_6)=(-U_6\bmod3)
\]
et la partition publiée des six coordonnées en deux triades.

Pour un mot \(w\), soient \(f(w)=|q^{-1}(w)|\) le nombre d'émissions,
\(m(w)\) sa masse de germes et
\[
 \nu(w)=(n_0(w),n_1(w),n_2(w))
\]
son recensement ordonné de trits. Soit en outre \(d(e)\) la classe diagonale additive d'une
émission et \(o(e)\in\{\mathrm{ES},\mathrm{NO}\}\) son orientation
publiée.

\begin{theorem}[Sélection orbitale interne des trois mots]
\label{thm:selector-orbital-constantes}
L'action \(D_3^{\rm cat}\) préserve l'ensemble des \(468\) émissions,
la projection \(q:E_{468}\to W_{243}\) et la multiplicité des germes. Son
action sur \(W_{243}\) possède exactement \(43\) orbites : \(38\) de taille
six et cinq de taille trois. Dans ce quotient, les unicités suivantes
sont vérifiées.

\begin{enumerate}
\item Il existe une unique orbite \(\mathcal O_\pi\) de taille six telle que, pour
tout \(w\in\mathcal O_\pi\),
\[
 f(w)=5,
 \qquad m(w)=1008,
 \qquad
 \{m(e):e\in q^{-1}(w)\}=\{432,144,144,144,144\}.
\]
C'est
\[
\mathcal O_\pi=
\{001112,010211,100121,112001,121100,211010\},
\]
et tous ses mots ont le recensement \(\nu=(2,3,1)\).

\item Considérons les orbites de taille six dont les mots ont une seule
émission, de multiplicité \(144\), et dont le support diagonal conjoint est
\(\{0,3,6\}\). Il y en a exactement six. Parmi elles, il existe une unique orbite
avec le même recensement \((2,3,1)\) que \(\mathcal O_\pi\) :
\[
\mathcal O_e=
\{011120,012110,101201,110012,120011,201101\}.
\]

\item Dans le même secteur, il existe une unique orbite équilibrée, c'est-à-dire avec
\(\nu=(2,2,2)\) :
\[
\mathcal O_\varphi=
\{002112,020211,112002,121200,200121,211020\}.
\]
\end{enumerate}

Enfin, la clause d'orientation
\begin{equation}
 \exists e\in q^{-1}(w):
 \qquad d(e)=6,
 \qquad o(e)=\mathrm{ES}
 \label{eq:gauge-orbital-constantes}
\end{equation}
sélectionne un unique représentant de chaque orbite et produit, respectivement,
\[
 \boxed{010211,\qquad201101,\qquad121200.}
\]
\end{theorem}

\begin{proof}
Les relations de groupe se vérifient directement comme des identités de
permutations de six coordonnées. L'énumération des \(468\) lignes
vérifie que \(r\) et \(s\) envoient chaque \(U_6\) sur une autre ligne du catalogue de
même multiplicité, et que la projection ternaire commute avec les deux
actions. Par conséquent, les orbites peuvent se calculer sur les \(243\) mots
sans perdre les fibres.

La partition exhaustive donne l'histogramme des tailles
\[
 38\cdot6+5\cdot3=243.
\]
En filtrant ces \(43\) orbites par le premier prédicat, il en reste exactement une ;
en filtrant par fibre à un point, multiplicité minimale et support diagonal
\(\{0,3,6\}\), il en reste six. Le recensement \((2,3,1)\) ne laisse qu'une orbite dans le
premier cas et le recensement équilibré \((2,2,2)\) n'en laisse qu'une dans le second.
Enfin, la recherche d'un relèvement avec \(d=6\) et orientation
\(\mathrm{ES}\) laisse un mot dans chacune des trois orbites. Toutes les
filtrations sont réalisées sur des descripteurs structurels du catalogue avant de consulter le
lecteur décimal.
\end{proof}

Les prédicats du théorème peuvent être réunis en un invariant de sélection.
Pour une orbite \(\mathcal O\), nous définissons son \emph{profil orbital}
\[
 \Sigma_{\rm cat}(\mathcal O)=
 \bigl(
 |\mathcal O|,\ f,\ m,\ \mathcal M_{\rm fib},\
 \mathcal D,\ \nu
 \bigr),
\]
où \(f\) est la taille de la fibre, \(m\) la masse des germes,
\(\mathcal M_{\rm fib}\) le multiensemble des multiplicités de ses
relèvements, \(\mathcal D\) le support diagonal lorsqu'il est employé et \(\nu\)
le recensement des trits. Les trois classes distinguées occupent des positions différentes
dans cet espace d'invariants :
\[
\begin{array}{c|c|c|c|c|c}
\mathcal O&|\mathcal O|&f&m&\mathcal M_{\rm fib}&(\mathcal D,\nu)\\ \hline
\mathcal O_\pi&6&5&1008&\{432,144,144,144,144\}&(-,(2,3,1))\\
\mathcal O_e&6&1&144&\{144\}&(\{0,3,6\},(2,3,1))\\
\mathcal O_\varphi&6&1&144&\{144\}&(\{0,3,6\},(2,2,2)).
\end{array}
\]
La singularité démontrée ne signifie pas que ces nombres reçoivent une
probabilité absolue sur \(\mathbb R\). Elle signifie qu'au sein de l'univers
fini produit par l'émetteur et avant l'évaluation des chiffres, leurs classes sont
séparables par des invariants internes et que la classe de \(\pi\) possède une fibre
régionale d'un type distinct. C'est la formulation mathématique d'une
priorité structurelle dans HMT\@.

Le contrôle négatif est informatif. Si l'on omet le support diagonal
\(\{0,3,6\}\), il existe quatre orbites à fibre unique équilibrées et
\(\mathcal O_\varphi\) cesse d'être unique. Avant de fixer la jauge, il y a six
orientations possibles dans chaque orbite. Après sa fixation, le mot de
\(\pi\) conserve trois relèvements positifs de \(U_6\), et non une seule région.
Ainsi, le théorème sélectionne le mot orienté mais respecte la multiplicité
géométrique des cinq régions.

La non-uniformité se renforce lors du passage aux orbites :
\[
\begin{array}{c|cc}
&\text{émissions}&\text{germes}\\ \hline
\mathcal O_\pi&5/78&14/243\\
\mathcal O_e&1/78&2/243\\
\mathcal O_\varphi&1/78&2/243.
\end{array}
\]
Ces quotients restent des masses du catalogue fini, et non des probabilités
attribuées aux nombres réels évalués.

L'avancée causale est précise. Les trois mots cessent d'être des entrées
nommées : ils sont retrouvés à partir de la symétrie et des invariants finis
d'APP--TPK, relativement à l'action \(D_3^{\rm cat}\), à l'étiquetage TRIT,
à la répartition \(3+3\) et à la jauge d'orientation de
\cref{eq:gauge-orbital-constantes}. La preuve n'utilise ni \(L_0\), ni \(L_1\),
ni chiffres décimaux, ni \(\alpha\), ni CODATA, ni l'état dodécaphasique. Restent des
affirmations ultérieures l'unicité absolue de cette action parmi tous les
automorphismes possibles, la déduction de la jauge au lieu de sa déclaration, la
canonicité du lecteur \(L_0/L_1\) indépendante de la valeur cible et l'émission coinductive des
développements infinis.

\section{Position de \texorpdfstring{\(\pi,e,\varphi\) et \(\alpha\)}{pi, e, phi et alpha} dans la classification}

Le \cref{thm:shift-hensel-completo} explique pourquoi tout préfixe
décimal peut être réémis à partir d'un état 3-adique : la capacité brute est
universelle. Par conséquent, la simple opération d'aller-retour sur un ruban ne sélectionne pas une
constante. Le contenu distinctif de la construction publiée se trouve
dans les objets antérieurs et postérieurs à cette réémission :
\begin{enumerate}
\item les fibres finies et les cinq régions distinctes de \(\pi\) ;
\item le lecteur archimédien typé qui associe mots, cylindres et blocs ;
\item la dodécaphase signée et la clôture réversible
\[
U_{12}^{\rm sgn}
\longleftrightarrow K
\longleftrightarrow(D_3K,D_4K,Q)
\longleftrightarrow\alpha_{\rm HMT};
\]
\item les critères de comparaison qui séparent une équivalence de cartes d'un
sélecteur antérieur à l'évaluation.
\end{enumerate}

Les douze triplets décimaux publiés de \(\pi,e,\varphi\) sont des coordonnées
finies exactes du lecteur enrichi. La chaîne imprimée de trente-six
chiffres est la projection dodécaphasique rationnelle
\(\alpha_{12}=\pi_{12}(\alpha_{\rm HMT})\) :
\[
\frac{7297352569283800997285105472380663}{10^{36}}.
\]
L'identité exacte dans cette fenêtre est
\[
\alpha_{12}
:=\pi_{12}(\alpha_{\rm HMT})
=\frac{7297352569283800997285105472380663}{10^{36}}.
\]
La récurrence compatible à profondeur arbitraire construit la section
infinie
\[
 \alpha_{\rm HMT}:=\varprojlim_N A^{(N)},
\]
et chaque fenêtre rationnelle est une projection exacte de ce même objet. Le
noyau \(E_{21/22}\) construit, dans une autre carte, la racine positive simple
finie \(\alpha_E\) ; le transport gradué du TPK prolonge ce noyau en
une famille compatible dont la racine limite est \(\alpha_B\). Si
\(\alpha_A:=\varprojlim_N\rho_N(\alpha_{12})\), les carrés de troncature
des deux familles commutent et déterminent à chaque profondeur le même
cylindre survivant. Le théorème exact des deux voies établit donc
\[
 \boxed{\alpha_A=\alpha_B=\alpha_{\rm HMT}}.
\]
L'égalité de \(\operatorname{Tr}_9(\alpha_E^{-1})\) et de
\(\operatorname{Tr}_9(\alpha_{12}^{-1})\) est le premier certificat fini de
cette identité ; elle ne la définit pas, ne limite pas sa portée et ne sélectionne pas les
prolongements. La projection typée conserve la distinction entre la section
coinductive, sa fenêtre rationnelle et la racine analytique sans séparer leur identité
à la limite.

La racine carrée de \(-1\) appartient à une autre carte. Dans le
\cref{chap:algebras-cuadraticas-orden-nueve}, on construit
\[
\mathbb F_9=\mathbb F_3[\iota]/(\iota^2+1),
\qquad \iota^2=-1,
\]
et l'on démontre que la matrice de Witt réalise cette structure quadratique. Il s'agit
d'une extension algébrique finie exacte, et non d'une catégorie visible de
développements réels.

\section{Carte de Hensel non filtrée, survie et frontière globale}

Les contrôles suivants sont des falsificateurs des résultats finis
précédents :
\begin{enumerate}
\item que deux états distincts modulo \(3^T\) émettent le même mot de
\(T\) blocs, ou qu'il existe un mot sans préimage ;
\item que l'inverse explicite de
\cref{eq:inversa-hensel-finita} échoue ;
\item qu'il apparaisse un rationnel dont le développement canonique ne soit pas ultimement
périodique, ou un irrationnel dont le développement le soit ;
\item qu'un émetteur déterministe fini produise un ruban non ultimement
périodique ;
\item que les recensements du catalogue ne totalisent pas \(468\), \(243\) et \(104\,976\),
ou que les mesures des trois mots ne coïncident pas avec l'énumération.
\end{enumerate}

La carte de Hensel brute
\[
\mathbb Z_3^6\longrightarrow[0,1]
\]
est continue et surjective ; en ajoutant la partie entière, son atlas brut couvre
\(\mathbb R\). Ce résultat exact n'implique pas que toute fibre brute contienne
un représentant qui survive aux filtres supplémentaires. Les fenêtres
finies publiées et les lecteurs nommés de \(\pi,e,\varphi\) sont vérifiés
dans HMT, relativement à leurs germes, régions,
orientations et règles enrichies.

\ifdefined\HMTInternalTraceability
Pour un préfixe filtré, soit \(T_u\) l'arbre des prolongements compatibles.
Si \(T_u\) est à ramification finie et possède un nœud à toute profondeur, le
lemme de König produit une branche infinie ; si, en outre, chaque niveau stable est
unitaire, la branche est unique. Cet énoncé est conditionnel : la non-vacuité universelle de \(T_u\)
pour tout préfixe n'a pas été démontrée. Restent donc
comme \textsc{programme ouvert} la couverture de toutes les fibres par des
survivants filtrés et la descente bien définie de la somme et du produit
natifs au quotient archimédien.
\fi
